% Values: PRECISE, FAITHFUL
% Evidence planned:
% - PRECISE: every alternative of the five bare grammar figures of
%   Chapters 2 and 3 (Expressions, Matches, Declarations and Bindings;
%   Patterns; Type Expressions; Modules, Headers, Imports and Exposing Lists;
%   Top-level Declarations and Bindings) appears here in the same order with
%   the same right-hand label; every row of the derived forms tables of
%   Appendix A (infixed application, conditional, conditional chain, &&, ||,
%   lambda with several patterns, negation, list expression, field accessor,
%   field selection, list pattern, function binding, let with several
%   definitions, module body as written) appears as an
%   alternative of its phrase class, marked with a dagger (B6).
% - FAITHFUL: precedence, layout and restrictions kept by reference to
%   Sections 2.x and 3.x, not copied, so that the two cannot disagree; the
%   abstract symbols ; and `and' (A17) noted as written by layout.

\chapter{Full Grammar}\label{ch:appendix-full-grammar}

The syntax of programs is given in Chapter~\ref{ch:programs}. The syntax of
long identifiers is that of Chapter~\ref{ch:syntax-core}.

The full grammar of the Core and of Modules consists of the grammar of
Chapters~\ref{ch:syntax-core} and~\ref{ch:syntax-modules}, in which each
derived form of Appendix~\ref{ch:appendix-derived-forms} is added as an
alternative of the phrase class it belongs to. The grammatical rules are shown
in Figures~\ref{fig:appendix-full-grammar-expressions},
\ref{fig:appendix-full-grammar-patterns},
\ref{fig:appendix-full-grammar-types},
\ref{fig:appendix-full-grammar-modules}
and~\ref{fig:appendix-full-grammar-top-level}. An alternative whose label
carries the mark \fgderived{} is a derived form, whose equivalent form is
given in Appendix~\ref{ch:appendix-derived-forms}; every other alternative is
a phrase of the bare language, and has the same label as in
Chapter~\ref{ch:syntax-core} or~\ref{ch:syntax-modules}. % B6
No subclass of \cls{Exp} is introduced: which expressions may be the operands
of an operator chain is stated in Section~\ref{sec:syntax-core-infix}.

The grammatical conventions are those of Chapter~\ref{ch:syntax-core}, namely:
\begin{itemize}
\item The brackets $\opt{\ }$ enclose optional phrases.
\item Alternative forms for each phrase class are in order of decreasing
  precedence. L (resp.\ R) means left (resp.\ right) association.
\item The infixed application, conjunction and disjunction alternatives are
  grouped by the fixities of Figure~\ref{fig:core-fixities}, as
  Section~\ref{sec:syntax-core-infix} describes, and not by their order in
  Figure~\ref{fig:appendix-full-grammar-expressions}.
\item Which item \kw{-} is a negation symbol, and which \kw{.} begins a field
  selection or a field accessor, is decided by lexical analysis
  (Section~\ref{sec:syntax-core-lexical}).
\item The precedence of types and of patterns, and the extent of lambda
  expressions, case expressions and \kw{let} expressions, are as stated with
  the grammar of the Core (Figure~\ref{fig:core-grammar-expressions} and the
  conventions before it); the extent of each phrase is decided by the layout
  rule (Sections~\ref{sec:syntax-core-layout}
  and~\ref{sec:syntax-modules-layout}).
\item The separator \kw{;} between match rules and between top-level
  declarations, and the connective \kw{and} between the members of a value
  binding or of a custom-type binding, are abstract syntax. Neither is an
  item of source text: match rules, the definitions of a \kw{let} and the
  declarations of a module are separated by layout, and \kw{and} is not a
  reserved word. Likewise an annotated binding
  $\var\ \kw{:}\ \ty\ \kw{=}\ \expm$ is written in source text as the
  annotation $\var\ \kw{:}\ \ty$ followed by a definition of \var. The
  derived alternatives $\kw{let}\ \valbind_1\ \cdots\ \valbind_n\ \kw{in}\
  \expm$ and $\topdec_1\ \cdots\ \topdec_n$ are the forms in which several
  definitions and declarations are written; in them each $\valbind_i$ is a
  definition, that is, a value binding without \kw{and}, and each
  $\topdec_i$ is a single declaration, as listed in
  Section~\ref{sec:appendix-derived-forms-modules}.
  % A17
\end{itemize}

The syntactic restrictions are those of
Sections~\ref{sec:syntax-core-restrictions}
and~\ref{sec:syntax-modules-restrictions}. They apply to phrases of the full
grammar, that is, to a program as written, before any derived form is
replaced by its equivalent form.

\begin{figure}[tbp]
\centering
\begin{grammar}
\gprod{\atexp}{\scon}{special constant}
\galt{\longvar}{value variable}
\galt{\longcon}{value constructor}
\galt{\kw{(}\opv\kw{)}}{operator value}
\galt{\kw{()}}{unit}
\galt{\kw{(}\ \expm_1\ \kw{,}\ \expm_2\ \opt{\kw{,}\ \expm_3}\ \kw{)}}{tuple}
\galt{\kw{[}\ \kw{]}}{empty list}
\galt{\kw{[}\expm_1\kw{,}\ \cdots\kw{,}\ \expm_n\kw{]}}{list ($n \geq 1$)\fgderived}
\galt{\kw{\kwlb}\ \opt{\exprow}\ \kw{\kwrb}}{record}
\galt{\kw{\kwlb}\ \var\ \kw{|}\ \exprow\ \kw{\kwrb}}{record update}
\galt{\kw{.}\lab}{field accessor\fgderived}
\galt{\atexp\kw{.}\lab}{field selection\fgderived}
\galt{\kw{-}\atexp}{negation\fgderived}
\galt{\kw{let}\ \dec\ \kw{in}\ \expm}{local declaration}
\galt{\kw{let}\ \valbind_1\ \cdots\ \valbind_n}{definitions\fgderived}
\galt{\quad\kw{in}\ \expm}{($n \geq 2$)}
\galt{\kw{(}\ \expm\ \kw{)}}{parentheses}
\noalign{\smallskip}
\gprod{\exprow}{\lab\ \kw{=}\ \expm\ \opt{\kw{,}\ \exprow}}{expression row}
\noalign{\smallskip}
\gprod{\expm}{\atexp}{atomic}
\galt{\expm\ \atexp}{application (L)}
\galt{\expm_1\ \opv\ \expm_2}{infixed application\fgderived}
\galt{\expm_1\ \kw{\&\&}\ \expm_2}{conjunction\fgderived}
\galt{\expm_1\ \kw{||}\ \expm_2}{disjunction\fgderived}
\galt{\kw{\textbackslash}\,\atpat\ \kw{->}\ \expm}{function}
\galt{\kw{\textbackslash}\atpat_1\cdots\atpat_n\ \kw{->}\ \expm}{function ($n \geq 2$)\fgderived}
\galt{\kw{if}\ \expm_1\ \kw{then}\ \expm_2\ \kw{else}\ \expm_3}{conditional\fgderived}
\galt{\kw{if}\ \expm_1\ \kw{then}\ \expm'_1}{conditional chain\fgderived}
\galt{\quad\kw{else}\ \kw{if}\ \expm_2\ \kw{then}\ \expm'_2\ \cdots}{($n \geq 2$)}
\galt{\quad\kw{else}\ \kw{if}\ \expm_n\ \kw{then}\ \expm'_n}{}
\galt{\quad\kw{else}\ \expm}{}
\galt{\kw{case}\ \expm\ \kw{of}\ \match}{case analysis}
\noalign{\smallskip}
\gprod{\match}{\mrule\ \opt{\kw{;}\ \match}}{match}
\noalign{\smallskip}
\gprod{\mrule}{\pat\ \kw{->}\ \expm}{match rule}
\noalign{\smallskip}
\gprod{\dec}{\valbind}{value declaration}
\noalign{\smallskip}
\gprod{\valbind}{\var\ \opt{\kw{:}\ \ty}\ \kw{=}\ \expm\ \opt{\kw{and}\ \valbind}}{variable binding}
\galt{\var\ \opt{\kw{:}\ \ty}\ \atpat_1\cdots\atpat_n}{function binding\fgderived}
\galt{\qquad\kw{=}\ \expm\ \opt{\kw{and}\ \valbind}}{($n \geq 1$)}
\galt{\atpat\ \kw{=}\ \expm}{destructuring binding}
\end{grammar}
\medskip
{\small\fgderived{} derived form, Appendix~\ref{ch:appendix-derived-forms}.\par}
\caption{Full Grammar: Expressions, Matches, Declarations and Bindings}
\label{fig:appendix-full-grammar-expressions}
\end{figure}

In the infixed application alternative, \opv{} is neither \kw{\&\&} nor
\kw{||}. The alternative \kw{-}\atexp{} is atomic, so that a negated
expression may be the argument of an application.

\begin{figure}[tbp]
\centering
\begin{grammar}
\gprod{\atpat}{\kw{\_}}{wildcard}
\galt{\scon}{special constant}
\galt{\var}{variable}
\galt{\longcon}{constructor}
\galt{\kw{()}}{unit}
\galt{\kw{(}\ \pat_1\ \kw{,}\ \pat_2\ \opt{\kw{,}\ \pat_3}\ \kw{)}}{tuple}
\galt{\kw{[}\ \kw{]}}{empty list}
\galt{\kw{[}\pat_1\kw{,}\ \cdots\kw{,}\ \pat_n\kw{]}}{list ($n \geq 1$)\fgderived}
\galt{\kw{\kwlb}\ \opt{\patrow}\ \kw{\kwrb}}{record}
\galt{\kw{(}\ \pat\ \kw{)}}{parentheses}
\noalign{\smallskip}
\gprod{\patrow}{\lab\ \opt{\kw{,}\ \patrow}}{pattern row}
\noalign{\smallskip}
\gprod{\pat}{\atpat}{atomic}
\galt{\longcon\ \atpat_1\ \cdots\ \atpat_n}{constructed pattern ($n \geq 1$)}
\galt{\pat_1\ \kw{::}\ \pat_2}{list construction (R)}
\galt{\pat\ \kw{as}\ \var}{layered}
\end{grammar}
\medskip
{\small\fgderived{} derived form, Appendix~\ref{ch:appendix-derived-forms}.\par}
\caption{Full Grammar: Patterns}
\label{fig:appendix-full-grammar-patterns}
\end{figure}

\begin{figure}[tbp]
\centering
\begin{grammar}
\gprod{\ty}{\tyvar}{type variable}
\galt{\kw{\kwlb}\ \opt{\tyrow}\ \kw{\kwrb}}{record type expression}
\galt{\kw{\kwlb}\ \tyvar\ \kw{|}\ \tyrow\ \kw{\kwrb}}{extensible record type expression}
\galt{\kw{()}}{unit type expression}
\galt{\kw{(}\ \ty_1\ \kw{,}\ \ty_2\ \opt{\kw{,}\ \ty_3}\ \kw{)}}{tuple type expression}
\galt{\longtycon\ \ty_1\ \cdots\ \ty_n}{type construction ($n \geq 0$)}
\galt{\ty\ \kw{->}\ \ty'}{function type expression (R)}
\galt{\kw{(}\ \ty\ \kw{)}}{parentheses}
\noalign{\smallskip}
\gprod{\tyrow}{\lab\ \kw{:}\ \ty\ \opt{\kw{,}\ \tyrow}}{type-expression row}
\end{grammar}
\caption{Full Grammar: Type Expressions}
\label{fig:appendix-full-grammar-types}
\end{figure}

Type expressions have no derived forms; Figure~\ref{fig:appendix-full-grammar-types}
is Figure~\ref{fig:core-grammar-types} unchanged. Modules, headers, imports and
exposing lists have no derived forms either: every module has a header, and
Figure~\ref{fig:appendix-full-grammar-modules} is
Figure~\ref{fig:grammar-modules} unchanged.
% A16, A30

\begin{figure}[tbp]
\centering
\begin{grammar}
\gprod{\modm}{\header\ \imports\ \topdec}{module}
\gprod{\header}{\kw{module}\ \modid\ \kw{exposing}\ \explist}{plain}
\galt{\kw{port}\ \kw{module}\ \modid\ \kw{exposing}\ \explist}{port module}
\gprod{\imports}{\kw{import}\ \modid\ \opt{\kw{as}\ \modid'}}{import}
\galt{\quad\opt{\kw{exposing}\ \explist}\ \imports}{}
\galt{}{empty}
\gprod{\explist}{\kw{(..)}}{everything}
\galt{\kw{(}\,\expitem_1\ \kw{,}\ \cdots\ \kw{,}\ \expitem_n\,\kw{)}}{explicit}
\gnote{$(n \geq 1)$}
\gprod{\expitem}{\var}{value}
\galt{\kw{(}\opv\kw{)}}{operator}
\galt{\tycon}{type}
\galt{\tycon\,\kw{(..)}}{open type}
\end{grammar}
\caption{Full Grammar: Modules, Headers, Imports and Exposing Lists}
\label{fig:appendix-full-grammar-modules}
\end{figure}

\begin{figure}[tbp]
\centering
\begin{grammar}
\gprod{\topdec}{\valbind}{value}
\galt{\kw{type}\ \datbind}{custom type}
\galt{\kw{type}\ \kw{alias}\ \typbind}{type alias}
\galt{\kw{port}\ \var\ \kw{:}\ \ty}{port}
\galt{\topdec_1\ \kw{;}\ \topdec_2}{sequential}
\galt{\topdec_1\ \cdots\ \topdec_n}{module body ($n \geq 2$)\fgderived}
\gprod{\datbind}{\tycon\ \tyvarseq\ \kw{=}\ \conbind}{custom-type binding}
\galt{\qquad\opt{\kw{and}\ \datbind}}{}
\gprod{\conbind}{\con\ \ty_1\ \cdots\ \ty_n\ \opt{\kw{|}\ \conbind}}{constructor binding}
\gnote{$(n \geq 0)$}
\gprod{\typbind}{\tycon\ \tyvarseq\ \kw{=}\ \ty}{type-alias binding}
\gprod{\tyvarseq}{\tyvar_1\ \cdots\ \tyvar_k}{type parameters}
\gnote{$(k \geq 0)$}
\end{grammar}
\medskip
{\small\fgderived{} derived form, Appendix~\ref{ch:appendix-derived-forms}.\par}
\caption{Full Grammar: Top-level Declarations and Bindings}
\label{fig:appendix-full-grammar-top-level}
\end{figure}

A top-level value declaration as written is a variable binding or a function
binding without \kw{and}; the function binding alternative of
Figure~\ref{fig:appendix-full-grammar-expressions} applies to it.
